\begin{abstract}
Medical language models are fluent but hard to audit. Knowledge graphs, combined with these language models, offer structured, source-grounded facts. However, recent work often treats a knowledge graph as proof of grounding without verifying whether relevant subgraphs were retrieved, used, or clinically safe. Our survey follows a small set of representative systems through their full lifecycle: graph construction and updating, indexing, retrieval, reasoning, verification, and post-deployment monitoring. We classify integration methods by where and how tightly graphs and models combine, distinguishing simple prompt-pasting from graph-aware encoders, and map the underlying ontologies, clinical datasets, and benchmarks supporting retrieval and tool-using agents. We find that graph benefits are highly conditional. Graphs help when filling genuine knowledge gaps, entity linking works, retrieval returns small relevant subgraphs, and evaluations use matched text-retrieval baselines. They help less, and sometimes may hurt, when coverage is thin, mappings fail, or retrieved context overpowers decisive evidence. Finally, we cover in detail practical failure modes, including weak citations, poor temporal handling, and closed-loop feedback, in which model proposals are written back into the graph and later retrieved as independent evidence.
\end{abstract}

\begin{IEEEkeywords}
Biomedical knowledge graph, large language model, graph retrieval-augmented generation, agentic systems, clinical artificial intelligence, evaluation, provenance.
\end{IEEEkeywords}
